% !PRINTED 625
\weberpage{625}
%<<CONT
dann und nur dann verschwinde, wenn unter ihren irreduziblen Faktoren die Funktion $F(\theta,z)$ vorkommt. Dies ist die Definition des Verschwindens, nach der, wenn man ganz korrekt sein wollte, eine jede Gleichung $G(\theta,z) = 0$ als Kongruenz nach dem Modul $F$ aufzufassen wäre. Um nicht weitläufig zu sein, wollen wir aber diese Ausdrucksweise hier nicht brauchen.

Wenn wir die Gleichung (1) durch $a_0$ dividieren, so erhält sie die Form
\weq{3}{f(\theta,z) = \theta^n + b_1\theta^{n-1} + b_2\theta^{n-2} + \cdots + b_{n-1}\theta + b_n,}
worin die $b_1, b_2, \ldots, b_n$ ganze oder gebrochene rationale Funktionen von $z$ sind.

Das System aller ganzen und gebrochenen rationalen Funktionen $\Phi(\theta,z)$ von $\theta$ und $z$, in denen der Nenner nicht durch $F$ teilbar ist, also nicht verschwindet, hat die Eigenschaft, sich durch Addition, Subtraktion, Multiplikation und Division (außer durch 0) zu reproduzieren, und bildet daher einen \gesperrt{Körper algebraischer Funktionen} (Bd.\,I, \S\,146), den ich mit $\Omega$ bezeichnen will.

Der Grad $n$ der irreduziblen Gleichung (2) oder (3), also der Gleichung niedrigsten Grades, der $\theta$ genügt, heißt der \gesperrt{Grad des algebraischen Körpers in bezug auf} $\theta$ oder auch der \gesperrt{Grad von} $\theta$.

Ist $\varphi(\theta)$ eine ganze Funktion von $\theta$, deren Koeffizienten ganze oder gebrochene rationale Funktionen von $z$ sind, so kann man durch Division eine Gleichung bilden:
\[
\varphi(\theta) = f(\theta)q(\theta) + r(\theta),
\]
worin $q(\theta)$ und $r(\theta)$ ebensolche Funktionen wie $\varphi$ sind, von denen jedoch $r(\theta)$ \gesperrt{höchstens vom Grade} $n-1$ ist. Wegen (3) ist dann
\weq{4}{\varphi(\theta) = r(\theta).}

Ist $\varphi(\theta)$ nicht durch $f(\theta)$ teilbar, so haben die beiden Funktionen keinen rationalen Teiler gemein, und man kann zwei Funktionen des Körpers $f_1(\theta), \varphi_1(\theta)$ so bestimmen, daß
\[
f(\theta)f_1(\theta) + \varphi(\theta)\varphi_1(\theta) = 1
\]
wird (Bd.\,I, \S\,6). Da nun $f(\theta) = 0$ ist, so folgt hieraus:
\weq{5}{\varphi_1(\theta) = \frac{1}{\varphi(\theta)}.}
